\begin{statementbox}
\begin{center}\begin{tikzpicture}[font=\small,>=stealth,line join=round]
\draw[gray!40,rounded corners=3pt] (0,0) rectangle (14.2,8.3);\node at (7.1,7.700000000000001) {隔2变为相反高度，隔4才回到原高度};
\begin{scope}[shift={(1.7153846153846155,3.8842857142857143)},x=0.7692307692307693cm,y=1.3142857142857143cm]\draw[->,gray] (-0.8,0)--(7,0) node[right]{$x$};\draw[->,gray] (0,-1.7)--(0,1.8) node[above]{$y$};\draw[blue,thick,solid,domain=-0.4:6.6,samples=180] plot(\x,{sin(deg(pi*\x/2))});\fill[orange!85!black] (1,1) circle(2.2pt);\fill[orange!85!black] (3,-1) circle(2.2pt);\fill[orange!85!black] (5,1) circle(2.2pt);\foreach \v/\h in {1/1,3/-1,5/1}{\draw[gray,dashed] (\v,0)--(\v,\h);\node[below=5pt] at (\v,0){$\v$};}\end{scope}\node[anchor=west,align=left,text width=5.0cm] at (8.4,4.0) {$f=\sin\frac{\pi x}2,\ a=2$\\[9pt]$f(1)=1$\\[9pt]$f(3)=-1$\\[9pt]$f(5)=1$\\[9pt]正周期可取4};
\node[text width=12.6cm,align=center] at (7.1,.65) {蓝线是满足条件的一条示例曲线；证明只用任意输入的两次变号。};
\end{tikzpicture}\end{center}


\textbf{这条结论解决什么问题？}

题目不给 $f(x)$ 的具体式子，只告诉我们“自变量增加固定的数，函数值会变成相反数”。这时仍然能判断函数值何时重新出现，还能把 $f(2026)$ 这样的大数字变成小数字计算。

\medskip
\textbf{已知条件}

设 $f(x)$ 在全体实数上都有定义，$a$ 是一个不等于 $0$ 的固定数，并且对\textbf{每一个实数 $x$}，都有
\[
\boxed{\ f(x+a)=-f(x)\ }.
\]

\textbf{核心结论：加两次 $a$，函数值回到原样。}
\[
\boxed{\ f(x+2a)=f(x)\quad(\forall x\in\mathbb R).\ }
\]
当 $a>0$ 时，$2a$ 是函数的一个正周期；当 $a<0$ 时，$2|a|$ 是一个正周期。所以可统一记成：
\[
\boxed{\ T=2|a|\ \text{是一个正周期（不一定最小）}.\ }
\]

\end{statementbox}\clearpage\begin{statementbox}\textbf{进一步：移动奇数次变号，偶数次还原。} 对任意整数 $n$，
\[
\boxed{\ f(x+na)=
\begin{cases}
 f(x),&n\text{ 为偶数},\\
 -f(x),&n\text{ 为奇数}.
\end{cases} }
\]
例如 $f(x+3a)=-f(x)$，$f(x+4a)=f(x)$。

\medskip
\textbf{不能顺手推出的东西}
\begin{itemize}[leftmargin=*,itemsep=0.25em]
\item 不能直接说 $2|a|$ 是\textbf{最小}正周期。
\item 不能只凭这个条件，就说图像关于 $\bigl(\frac a2,0\bigr)$ 中心对称。
\item 若 $f$ 不是恒为 $0$，$|a|$ \textbf{本身不是周期}，因为移动这么远会变号。
\end{itemize}
\end{statementbox}
\clearpage
\begin{statementbox}
\textbf{补充｜错位反号型：两个输入都带偏移}

设 $a,b$ 是固定实数，$a\ne b$。先取 $f$ 定义在 $\mathbb R$ 上，且对每个实数 $x$ 都有
\[f(x+a)=-f(x+b).\]
则
\[\boxed{T=2|a-b|\text{ 是一个正周期，不一定最小。}}\]
对于其他非空定义域 $D$，须保证：令 $d=a-b$ 后，每个 $t\in D$ 都能在题设中取 $x=t-b$，且所得 $f(t+d)=-f(t)$ 对每个 $t\in D$ 成立；同时 $D+d=D$，即来回移动 $d$ 都仍在域内。不能只在两边碰巧都有定义的局部使用等式，就宣布整个函数有周期。
\begin{center}\begin{tikzpicture}[font=\small,>=stealth,line join=round]
\draw[gray!40,rounded corners=3pt] (0,0) rectangle (14.2,8.3);
\node[text width=12.6cm,align=center] at (7.1,7.65) {错位反号：先认两个输入，再认两次移动};

\begin{scope}[shift={(1.8,4.3)},x=1.45cm,y=1.05cm]
\draw[->,gray] (-.4,0)--(6.8,0) node[right]{$u$};
\draw[->,gray] (0,-1.5)--(0,1.5) node[above left=3pt]{$y$};
\draw[blue,thick,domain=0:6.3,samples=160] plot(\x,{cos(deg(pi*\x/3))});
\foreach \u/\v/\lab in {0/1/t,3/-1/{t+d},6/1/{t+2d}}{
\draw[gray,dashed] (\u,0)--(\u,\v);\fill[orange!85!black] (\u,\v) circle(2.3pt);
\node[below=5pt] at (\u,-1.55) {$\lab$};}
\draw[<->,violet] (0,1.65)--(3,1.65) node[midway,above]{$d$};
\draw[<->,violet] (3,1.65)--(6,1.65) node[midway,above]{$d$};
\end{scope}
\node[align=center] at (7.1,1.3) {$t=x+b,\quad t+d=x+a,\quad d=a-b$};

\node[text width=12.6cm,align=center] at (7.1,.65) {图取 $d=3,t=0$，蓝线为 $y=\cos(\pi u/3)$；一般证明不依赖这条示例曲线。};
\end{tikzpicture}\end{center}


\end{statementbox}
